% ===== 原文 pp.84–87（PDF p.94–97） =====
% [原文 p.84]
\chapter{修正 A-模型}

本章详细研究把 Vafa-Witten 理论从 4d 紧化到 2d、紧化在亏格 $g_C\geqslant 2$ 的黎曼面 $C$ 上所得到的 2d TFT。我们称这个 2d TFT 为\emph{修正 A-模型}（modified A-model），理由是它可由靶空间为 Kähler 流形 $X$ 的普通 2d A-模型经 $X$ 的一个 $U(1)$ 等距作用扭转而得到。据我们所知，修正 A-模型在本文之外鲜有研究，因此我们详细讨论它。

注意，在附录 A 中我们给出另一系列 2d TFT，它们出现在对各扭转版本的 $d=4$、$\cN=4$ SYM 作维数约化的分析中（其中包括 A-型与 B-型 2d 拓扑规范理论，以及带规范场的 2d B-模型；后者已在上章的面算子分析中露过面）。

\para{扭转 2d A-模型}

我们已注意到 $t=0$ 的 GL-扭转理论与 Vafa-Witten 理论之间颇有相似之处。文献 [35] 证明：随着 $C$ 的尺寸收缩到零，乘积 $\Sigma\times C$ 上的 $t=0$ 理论紧化为 $\Sigma$ 上以 $\cM_H(G,C)$ 为靶空间的 A-模型，其中 $\cM_H(G,C)$ 是黎曼面 $C$ 上 Hitchin 方程的解的模空间。（空间 $\cM_H(G,C)$ 大致可看作 $C$ 上平坦 $G$-丛的模空间的余切丛；我们将在下一章进一步讨论它。）因此我们预期，对 Vafa-Witten 理论在 $C$ 上的紧化会得到类似的描述；事实也确实如此--只是带一点扭转。

我们得到的不是 A-模型，而是一个我们称为\emph{修正 A-模型}的理论，它是 A-模型的一个扭转版本。

我们从零开始构造以靶空间 $X$ 为目标的修正 2d A-模型，其中 $X$ 是任意 Kähler 流形，配备群 $U(1)$ 的一个 Hamilton 作用，并假定该作用同时保持 $X$ 的 Kähler 形式与复结构（从而也保持 Kähler 度规）。从 $\Sigma$ 上以 $X$ 为靶空间的 A-模型出发，通过把 $\Sigma$ 的旋转群 $U(1)_E$ 替换为

% [原文 p.85]
\begin{equation}
U(1)_E \to U(1)'_E = \diag\big(U(1)_E\times U(1)_X\big), \tag{7.1}
\end{equation}

来扭转理论；其中 $U(1)_X$ 是靶空间上 $U(1)$ 作用在场上诱导的变换，$\diag$ 表示对角子群。等价地，也可以从 $\Sigma$ 上以 $X$ 为靶空间的 (2,2) 非线性 $\sigma$-模型出发，作替换

\[
U(1)_E \to U(1)'_E = \diag\big(U(1)_E\times U(1)_V\times U(1)_X\big),
\]

其中 $U(1)_V$ 是未扭转理论的向量 R-对称性。

局部地看 $\Sigma$（即只关注 $\Sigma$ 的主 $U(1)$ 标架丛的单个平凡化 chart），修正 A-模型与普通 A-模型几乎完全相同。两个理论的差别只有在考虑 chart 之间的重叠时才显现出来：$\Sigma$ 上的旋转伴随有 $U(1)_X$ 在场上作用。

本章考虑的修正 A-模型与具一般规范丛和涨落规范场的带规范场 $\sigma$-模型（例如第 3 章考虑的带规范场 4d A-模型）是近亲。特别地，文献 [12] 等处讨论过带规范场的 2d A-模型。修正 2d A-模型是它的一个特化：规范丛取为 $\Sigma$ 上的主 $U(1)$ 标架丛 $P$，联络固定为（与度规相容、无挠的）自旋联络。

7.1 节回顾普通 2d A-模型。7.2 节详细描述修正 A-模型的扭转，并写出扭转理论的场含量、变分与作用量；我们将发现有必要用一项额外项来扩充 A-模型通常的拓扑项。7.3 节给出这个等变拓扑项的一个解释。7.4 节分析修正 A-模型的局部可观测量谱与路径积分测度的鬼数反常。随后 7.5 节考虑定义在带边界黎曼面上的修正 A-模型并讨论膜的范畴。最后 7.6 节分析当边界非平凡时鬼数对称性无反常的条件。

% [Tutor Note: p.85 式 (7.1) 的“对角”扭转要按场的自旋逐项落实：$z,\bar z$ 在 $U(1)_E$ 下的荷为 $(+1,-1)$，靶空间复坐标 $\varphi^i$ 在 $U(1)_X$ 下的荷为 $q_i$，扭转后场的新自旋为 $s'=s+\tfrac12 q$。两个可立刻自检的退化情形写在正文后面：$X$ 为单点时 $U(1)_X$ 平凡、(7.1) 退回普通 A-模型扭转；p.91 末指出 $\Sigma$ 为平坦环面时自旋联络平凡、作用量退回普通 A-模型。这两条都只能在对角化 (7.1) 的前提下成立。]
% [Tutor Note: p.85 (7.1) 下面那个无编号替换多出一个 $U(1)_V$：因为普通 A-模型本身就是把 (2,2) $\sigma$-模型按向量 R-对称性 $U(1)_V$ 扭转而来（见 p.87 开头）。两种起点的差别正是“先按 $U(1)_V$ 扭转、再对 $U(1)_E$ 对角化”与“一次对 $U(1)_E\times U(1)_V\times U(1)_X$ 对角化”的次序差别，最终给出同一个修正 A-模型。]

\section{2d A-模型}

我们先回顾熟悉的 2d A-模型的场含量、作用量与变分。设 $\Sigma$ 是亏格 $g_\Sigma$、紧致、可定向、无边界的黎曼面，$h$ 是 $\Sigma$ 上的黎曼度规。

% [原文 p.86]
设 $z,\bar z$ 是 $\Sigma$ 上适配 $h$ 的局部复坐标，即
\[
h_{zz}=h_{\bar z\bar z}=0,\qquad h_{z\bar z}=h_{\bar z z}=\tfrac12\ee^{2f},\qquad h^{z\bar z}=h^{\bar z z}=2\ee^{-2f},
\]
其中 $f=f(z,\bar z)$ 是局部实函数。记 $\sigma^1=\Re z$、$\sigma^2=\Im z$ 为 $\Sigma$ 上适配于 $z,\bar z$ 的实坐标。$h$ 的实分量为 $h_{\mu\nu}=\ee^{2f}\delta_{\mu\nu}$，$\mu,\nu=1,2$。$\Sigma$ 上的 Kähler（体积）形式为
\[
\varepsilon=\ii h_{z\bar z}\dd z\wedge\dd\bar z=\frac{\ii}{2}\ee^{2f}\dd z\wedge\dd\bar z=\ee^{2f}\dd\sigma^1\wedge\dd\sigma^2=\sqrt{h}\,\dd^2\sigma,
\]
最后一个表达式对 $\Sigma$ 上任意实坐标 $\sigma^\mu$ 都成立。

靶空间 $X$ 取为 Kähler 流形，其 Kähler 形式为
\[
\Omega=\ii g_{i\bar j}\dd\varphi^i\wedge\dd\varphi^{\bar j}=\frac12\Omega_{IJ}\dd\varphi^I\wedge\dd\varphi^J,
\]
其中 $\varphi^I$ 是 $X$ 上的局部实坐标，$\varphi^i,\varphi^{\bar j}$ 是 $X$ 上的局部复坐标，$g$ 是 Kähler 度规，满足
\[
g_{ij}=g_{\bar i\bar j}=0,\qquad g_{i\bar j}=g_{\bar j i}=(g_{j\bar i})^*=(g_{\bar i j})^* .
\]

A-模型的玻色场由一个映射
\[
\Phi:\Sigma\to X
\]
（局部由函数 $\varphi^I(\sigma)$ 描述）以及辅助场
\begin{gather*}
P\in\Omega^{0,1}\big(\Sigma,\Phi^*T^{1,0}X\big)\\
\tilde P\in\Omega^{1,0}\big(\Sigma,\Phi^*T^{0,1}X\big)
\end{gather*}
组成，其中 $T^{1,0}X$ 与 $T^{0,1}X$ 是 $X$ 的复化切丛关于 $X$ 上复结构 $J_X$ 的全纯与反全纯子丛；类似地，记号 $\Omega^{1,0}$ 与 $\Omega^{0,1}$ 表示 1-形式在 $\Sigma$ 上按复结构 $J_\Sigma$ 到 $(1,0)$ 与 $(0,1)$ 部分的投影。

A-模型的费米场由切丛取值的微分形式
\begin{gather*}
\chi\in\Omega^0\big(\Sigma,\Phi^*TX\big)\\
\rho\in\Omega^{0,1}\big(\Sigma,\Phi^*T^{1,0}X\big)\\
\tilde\rho\in\Omega^{1,0}\big(\Sigma,\Phi^*T^{0,1}X\big)
\end{gather*}
组成。

% [Tutor Note: p.86 第一组度规公式可这样自检：取 $z=\sigma^1+\ii\sigma^2$，则 $\dd z\wedge\dd\bar z=-2\ii\,\dd\sigma^1\wedge\dd\sigma^2$，于是 $\ii h_{z\bar z}\dd z\wedge\dd\bar z=\ii\cdot\tfrac12\ee^{2f}\cdot(-2\ii)\dd\sigma^1\wedge\dd\sigma^2=\ee^{2f}\dd\sigma^1\wedge\dd\sigma^2$，与 $\sqrt h=\ee^{2f}$（因 $\det h_{\mu\nu}=\ee^{4f}$）一致。若改用 $z=(\sigma^1+\ii\sigma^2)/\sqrt2$，$\sqrt h$ 前会出现因子 $2$，故作者此处取的必是前一种坐标归一。另可核对 $h^{z\bar z}=2\ee^{-2f}$ 恰为 $h_{z\bar z}=\tfrac12\ee^{2f}$ 的逆。]
% [Tutor Note: p.86 中 $\Omega=\tfrac12\Omega_{IJ}\dd\varphi^I\wedge\dd\varphi^J$ 的 $\tfrac12$ 是实的反对称约定。第 88 页用 $\partial_JH=\Omega_{IJ}V^I$ 时这个约定会自动“吃掉”因子：$\iota_V\Omega=\tfrac12\Omega_{IJ}(V^I\dd\varphi^J-V^J\dd\varphi^I)=\Omega_{IJ}V^I\dd\varphi^J$。由此可立刻反推：分量式 $\partial_JH=\Omega_{IJ}V^I$ 等价于 $\dd H=+\iota_V\Omega$，而 (7.5) 印作 $\dd H=-\iota_V\Omega$，两者相差一个符号（详见 p.88 批注）。]

% [原文 p.87]
由于 A-模型是把 (2,2) 超对称非线性 $\sigma$-模型按其 $U(1)_V$ 向量 R-对称性扭转得到的，$U(1)_A$ 轴 R-对称性在扭转理论中仍是对称性，只是在那里我们改称它为鬼数对称性。玻色场的鬼数为 0，费米场 $\chi$ 的鬼数为 1，而场 $\rho$ 与 $\tilde\rho$ 的鬼数为 $-1$。BRST 变分 $\delta_Q$ 本身携带一个单位的鬼数；它在场上作用为
\begin{equation}
\begin{aligned}
\delta_Q\varphi^i &= \chi^i, &
\delta_Q\varphi^{\bar i} &= \chi^{\bar i},\\
\delta_Q\chi^i &= 0, &
\delta_Q\chi^{\bar i} &= 0,\\
\delta_Q\rho_{\bar z}{}^i &= P_{\bar z}{}^i-\Gamma^i{}_{jk}\chi^j\rho_{\bar z}{}^k, &
\delta_Q\tilde\rho_z{}^{\bar i} &= \tilde P_z{}^{\bar i}-\Gamma^{\bar i}{}_{\bar j\bar k}\chi^{\bar j}\tilde\rho_z{}^{\bar k},\\
\delta_QP_{\bar z}{}^i &= -R^i{}_{jk\bar l}\rho_{\bar z}{}^j\chi^k\chi^{\bar l}-\Gamma^i{}_{jk}\chi^jP_{\bar z}{}^k, &
\delta_Q\tilde P_z{}^{\bar i} &= -R^{\bar i}{}_{\bar j\bar k l}\tilde\rho_z{}^{\bar j}\chi^{\bar k}\chi^l-\Gamma^{\bar i}{}_{\bar j\bar k}\chi^{\bar j}\tilde P_z{}^{\bar k},
\end{aligned} \tag{7.2}
\end{equation}
其中 $\Gamma^i{}_{jk}$、$\Gamma^{\bar i}{}_{\bar j\bar k}$、$R^i{}_{jk\bar l}$、$R^{\bar i}{}_{\bar j\bar k l}$ 是 $X$ 上 Kähler 度规的非零 Christoffel 符号与 Riemann 张量系数（看作 $\Phi$ 的函数）。容易验证，上述变分的离壳形式平方为零。

% [Tutor Note: p.87 (7.2) 的“离壳平方为零”是本章最需要动手核的一处。核对路线：$\delta_Q^2\rho_{\bar z}{}^i=\delta_QP_{\bar z}{}^i-(\partial_l\Gamma^i{}_{jk})\chi^l\chi^j\rho_{\bar z}{}^k-\Gamma^i{}_{jk}\chi^j\delta_Q\rho_{\bar z}{}^k$，其中第一项与第三项各产生一个含 $\chi P$ 的项，$\partial_l\Gamma$ 项则要靠 $\chi^l\chi^j$ 的反对称性与 Riemann 张量恒等式 $R^i{}_{jkl}=\partial_k\Gamma^i{}_{jl}-\partial_l\Gamma^i{}_{jk}+\Gamma\Gamma-\Gamma\Gamma$ 化成曲率项。按 (7.2) 印出的符号，两条 $\delta_QP$ 的 Christoffel 项都写作负号，读者若发现 $\chi P$ 项不能相消，应以“$\delta_Q^2=0$”为准反推该项符号（本译本照录原文）。]
% [Tutor Note: p.87 (7.3) 下方 $V$ 中系数 $4\ii$ 的来历可这样核对：$V$ 里出现的是 $\rho_{\bar z}{}^i(\tilde P_z{}^{\bar j}+4\ii\partial_z\varphi^{\bar j})$ 这一组合，代入 p.87 给的运动方程 $\tilde P_z{}^{\bar j}=-2\ii\partial_z\varphi^{\bar j}$ 后它化为 $2\ii\rho_{\bar z}{}^i\partial_z\varphi^{\bar j}$；也就是说 $4\ii$ 与 $2\ii$ 之差恰好把辅助场消去后残留的 $\rho\,\partial\varphi$ 耦合定在 $2\ii$ 上。独立的核对法：直接把 $\delta_QV$ 展开（用 (7.2) 的 $\delta_Q\rho$、$\delta_QP$），与 p.91 的展开式逐项比对 $P\tilde P$、$P\partial\varphi$ 与 $\rho\partial\chi$ 的系数，即可唯一确定这个 $4\ii$。]

A-模型的作用量是一个 BRST-正合项与一个与度规无关的项之和：
\begin{equation}
S=\delta_QV+t\int_\Sigma\Phi^*(\Omega) \tag{7.3}
\end{equation}
其中
\[
V=-\frac{1}{4t}\int_\Sigma g_{i\bar j}\Big\{\rho_{\bar z}{}^i\big(\tilde P_z{}^{\bar j}+4\ii\partial_z\varphi^{\bar j}\big)+\tilde\rho_z{}^{\bar j}\big(P_{\bar z}{}^i-4\ii\partial_{\bar z}\varphi^i\big)\Big\}(\ii\dd z\wedge\dd\bar z),
\]
而 $t$ 是乘在整个作用量前面的耦合参数。$\sigma$-模型映射下拉回的 Kähler 形式在局部坐标下由下式给出：
\begin{equation}
\int_\Sigma\Phi^*(\Omega)=\int_\Sigma\frac12\Omega_{IJ}\dd\varphi^I\wedge\dd\varphi^J=\int_\Sigma g_{i\bar j}\Big\{\partial_z\varphi^i\partial_{\bar z}\varphi^{\bar j}-\partial_{\bar z}\varphi^i\partial_z\varphi^{\bar j}\Big\}(\ii\dd z\wedge\dd\bar z). \tag{7.4}
\end{equation}
用运动方程
\[
P_{\bar z}{}^i=2\ii\partial_{\bar z}\varphi^i,\qquad \tilde P_z{}^{\bar i}=-2\ii\partial_z\varphi^{\bar i}
\]
消去辅助场后可见，A-模型的路径积分局域化在全纯瞬子方程
\[
\partial_{\bar z}\varphi^i=0
\]
的解上，后者对应于复流形之间的全纯映射 $\Phi:\Sigma\to X$。

% [Tutor Note: p.87 末尾把“全纯瞬子方程”与其解“全纯映射”等同，理由是把 $\partial_{\bar z}\varphi^i=0$ 看成 $\bar\partial\Phi=0$，即映射在复结构 $J_\Sigma,J_X$ 下满足 $\dd\Phi\circ J_\Sigma=J_X\circ\dd\Phi$（p.91 对修正模型写出的是同一条件的协变化版本）。核对时注意 $\varphi^i$ 只取全纯坐标分量，反全纯分量 $\varphi^{\bar i}$ 的方程是共轭的，不独立。]
